% 18 November 2026. Drafted for Edward Baker with OpenAI Codex assistance.
% Actual model/effort unavailable; helper settings inherited.
\section{Backward comparison and regional gap statistics}
\label{sec:statistics}

The preceding hierarchy describes the exterior of a specified pair. A
statistical question instead asks how the smallest gap in a region evolves,
allowing its minimizing pair to change. These viewpoints can be connected
without assigning randomness to the zeros: choose a region at random,
retain the complete deterministic evolution in each region, and average
its observables. The essential additional information is the environment
of a small gap. In the all-real regime the hierarchy gives that environment
a useful sign, leading to a comparison determined by three initial gaps.
We develop the local comparison before forming its statistical consumers,
following the conditional results of \cite{BakerCM16,BakerCM17}.

\subsection{What minimum selection does and does not provide}

The strong collar already used in this paper makes the target occurrences
strict mutual nearest neighbors. Indeed, under \eqref{eq:domain},
\begin{equation}
 |\rho_j-z_\pm|\ge r_0-|s|\ge\frac{3r_0}{4}
 >\frac{r_0}{2}\ge |z_+-z_-|.
 \label{stat:mutual-neighbors}
\end{equation}
This compares each endpoint with its exterior; it does not compare two
exterior zeros. Nor does it require an attained global minimum.

\begin{proposition}[Complete versus regional minima]
\label{stat:minimum-selection}
At a simple canonical slice in $0\le t\le1/2$, the infimum of distances
between distinct zero occurrences is zero and is not attained. A finite
set of at least two occurrences has a minimum; if its points are real and
ordered, a minimizing pair is consecutive in that finite ordering.
If a conjugation-invariant target $c\pm s$ is shortest among a specified
comparison set and $s\ne0$, every compared exterior occurrence satisfies
\begin{equation}
 |w_j-s|,|w_j+s|\ge2|s|,
 \qquad |Q_j|\ge4|D|,
 \qquad |w_j|\ge\sqrt{3|D|}.
 \label{stat:relative-collar}
\end{equation}
\end{proposition}

\begin{proof}
The complete canonical count on the inherited slab is
$N_t(X)\sim X\log X/(4\pi)$ for positive real parts, with all zeros in
$|\Im z|\le1$; the closed-slab count passage is given in
\cite{BakerGT33}. If all distinct zeros were separated by $g>0$, their
open disks of radius $g/2$ would be disjoint. Disks centered at zeros with
$0<\Re z\le X$ lie in a rectangle of area
$(X+g)(2+g)$. Comparison of areas would give $N_t(X)=O_g(X)$,
contrary to the count. Simplicity excludes an attained zero distance.
Existence of a finite minimum is immediate. For real ordered points,
a nonconsecutive distance is a sum of consecutive positive gaps and
cannot be smaller than all of them. Finally, shortest-pair selection
gives the first inequalities in \eqref{stat:relative-collar}; their
product gives the bound for $Q_j$. Adding their squares gives
$2|w_j|^2+2|s|^2\ge8|s|^2$, proving the last bound.
\end{proof}

The separation in \eqref{stat:relative-collar} vanishes with the target
gap. It does not supply the independent collar in \eqref{eq:domain},
which gives the stronger fixed estimate $|Q_j|\ge15r_0^2/16$.
For a window minimum, the comparison says nothing about unexamined
exterior occurrences. At a complex slice, a closest pair also need not be
conjugation invariant: the points $\ii h,\varepsilon+\ii h,-\ii h,
\varepsilon-\ii h$, with $0<\varepsilon<2h$, have horizontal closest
pairs. Even an admissible closest pair need not detect all nonreality;
$\{-\varepsilon,\varepsilon,-\ii,\ii\}$ with
$0<\varepsilon<1/2$ has a real closest pair and a nonreal pair.
These finite examples test selection alone and do not satisfy all the
canonical hypotheses.

For perspective, at an all-real ordered heat-flow slice satisfying the
pair law, with an attained positive global minimum gap $g=2s$, the $k$th exterior zero on either side
of that pair has center-distance at least $(2k+1)s$. Thus, with fewer
terms allowed on a finite side,
\begin{equation}
 0\le M\le2\sum_{k\ge1}\frac1{4k(k+1)D}=\frac1{2D},
 \qquad D_t=2-4DM\ge0.
 \label{stat:telescoping}
\end{equation}
This useful extremal calculation belongs to its stated class. The
canonical simple divisor has no such attained global minimum, and
forward nondecrease does not furnish a backward margin.

\subsection{Adjacent real pairs require only strict isolation}

We now change the pair domain explicitly. Fix
\begin{equation}
 0\le t_*-\Delta\le t_*\le\tfrac12,
 \qquad u=t_*-t.
 \label{stat:slab}
\end{equation}
Here $u$ is \emph{physical backward time}, with no factor $\ell^2$.
Suppose the complete divisor at $t_*$ is real and simple. Follow an
adjacent pair $x_j(u)<x_{j+1}(u)$ on a common interval $0\le u<\tau$
on which the complete divisor remains real and simple and the branches
continue. This is a stopped regime, not a conclusion of initial reality
at one time. Set
\begin{equation}
 \begin{gathered}
 c_j=\frac{x_j+x_{j+1}}2,\qquad
 D_j=\frac{(x_{j+1}-x_j)^2}{4},\\
 Q_{j,\rho}=(\rho-c_j)^2-D_j
             =(\rho-x_j)(\rho-x_{j+1}),\qquad
 M_j=\sum_{\rho\ne x_j,x_{j+1}}Q_{j,\rho}^{-1},\qquad
 S_{j,n}=\sum_{\rho\ne x_j,x_{j+1}}Q_{j,\rho}^{-n}.
 \end{gathered}
 \label{stat:real-moments}
\end{equation}
All sums include every exterior occurrence. The canonical product and
inverse-square tails are still those of the complete source.
We do not impose the factor-four collar or a prescribed height interval
on the adjacent pair in this subsection.

\begin{lemma}[Strict-isolation passage]
\label{stat:strict-isolation}
On every compact subinterval of the stopped regime,
\begin{equation}
 \frac{dD_j}{du}=-2+4D_jM_j,
 \qquad
 \frac{dM_j}{du}=-2M_j^2+12S_{j,2}+16D_jS_{j,3},
 \qquad M_j,S_{j,2},S_{j,3}\ge0.
 \label{stat:backward-system}
\end{equation}
These identities do not require the quantitative collar of
\eqref{eq:domain}.
\end{lemma}

\begin{proof}
At a reference time the pair is adjacent, so every exterior real zero
has $|\rho-c_j|>\sqrt{D_j}$. Local finiteness supplies a radius
\[
 \sqrt{D_j}<R<\min_{\rho\ne x_j,x_{j+1}}|\rho-c_j|.
\]
The circle of radius $R$ about this reference center contains exactly
the two target occurrences and is zero-free. Hold it fixed while
differentiating in time. The argument-principle moments
\eqref{eq:coordinates} define $c_j,D_j$ locally, and the deflated
logarithmic derivative gives
\[
 M_j=-\frac1{2\pi\ii}\oint_\Gamma
              \frac{H_z/H}{(z-c_j)^2-D_j}\dd z.
\]
To pass the finite polynomial identity to $H$, freeze this time and
truncate the complete canonical product in opposite and conjugate groups,
retaining the target and all its occurrences. Use each truncated product
as initial data for its exact polynomial heat evolution and differentiate
at the reference time. The products and all fixed spatial derivatives
converge locally uniformly. On $\Gamma$,
$H$ and the target denominator are bounded away from zero. Hence the
logarithmic derivatives and their time derivatives converge, using
\[
 \partial_t(H_z/H)=-H_{zzz}/H+H_zH_{zz}/H^2.
\]
The contour moments and their derivatives consequently converge as well.
The finitely many nearby exterior denominators are nonzero, while the
remaining summands are $O(|\rho|^{-2n})$ and have convergent complete
tails. Thus the finite identities for $D_t$ and $M_t$, obtained in
\eqref{eq:pairlaw} and \eqref{eq:low-orders}, pass to the limit.
Their polynomial derivation uses strict separation, not the factor-four
ratio. Reversing time gives \eqref{stat:backward-system}. Finally,
$Q_{j,\rho}>0$ because the exterior is real and lies outside the adjacent
pair; this proves the signs. Finitely many local charts cover any compact
pre-collision interval.
\end{proof}

Only the identities have been extended. None of the earlier quantitative
collar estimates is asserted on this larger domain. The polynomial
passage is made at a frozen time through spatial jets; it is not a
common-time convergence theorem for finite polynomial heat evolutions.
A uniform isolating circle need not survive to a multiple collision.
The comparisons below approach such an event by continuity of the real
gap, without differentiating its square-root labeling at the event.

\subsection{A backward comparison from initial exterior mass}

\begin{theorem}[Initial-environment comparison]
\label{stat:backward-comparison}
For one pair in the stopped real regime, suppose
$0<d_0\le D(0)$ and $0\le b_0\le M(0)$. Put
$a(u)=1+2b_0u$. Then
\begin{equation}
 M(u)\ge\frac{b_0}{1+2b_0u},\qquad
 \boxed{D(u)\ge a(u)\bigl[d_0a(u)-2u\bigr]}
 \quad(0\le u<\tau).
 \label{stat:comparison}
\end{equation}
If
\begin{equation}
 d_0(1+2b_0\Delta)>2\Delta,
 \label{stat:strict-margin}
\end{equation}
this pair cannot be a colliding pair at a first finite collision time
$u_c\le\Delta$ up to which the stipulated all-real regime holds.
\end{theorem}

\begin{proof}
Positivity in \eqref{stat:backward-system} gives $M_u\ge-2M^2$.
The function $b(u)=b_0/(1+2b_0u)$ satisfies $b_u=-2b^2$, whence
\[
 (M-b)_u+2(M+b)(M-b)\ge0.
\]
Multiplying by its positive integrating factor, on each compact
pre-collision interval, proves $M\ge b$. Since $D\ge0$ there,
\[
 (a^{-2}D)_u=a^{-2}(D_u-4bD)\ge-2a^{-2}.
\]
The exact integral $\int_0^u a(v)^{-2}\dd v=u/a(u)$, including $b_0=0$,
gives \eqref{stat:comparison}. The bracket
$d_0+(2b_0d_0-2)u$ is affine, positive at zero and, under
\eqref{stat:strict-margin}, positive at $\Delta$. It is therefore
strictly positive on the full interval. If the pair were to collide
at $u_c\le\Delta$, its continuous gap would give $D(u)\to0$ as
$u\uparrow u_c$, contradicting the strictly positive lower bound.
\end{proof}

When $b_0d_0<1$, the explicit sufficient duration is
\begin{equation}
 \Delta<\frac{d_0}{2(1-b_0d_0)}.
 \label{stat:duration}
\end{equation}
When $b_0d_0\ge1$, the scalar inequality places no finite cap on
$\Delta$; the stopped-regime and coverage hypotheses still apply.
In particular, the theorem does not continue a pair's positivity
argument past an earlier unprotected collision elsewhere.

The same calculation bounds the exact backward transport cost. With
$A(u)=\int_0^u M(v)\dd v$, integration of the first equation in
\eqref{stat:backward-system} yields
\begin{equation}
 \begin{split}
 D(u)&=\e^{4A(u)}
       \left[D(0)-2\int_0^u\e^{-4A(v)}\dd v\right],\\
 2\int_0^u\e^{-4A(v)}\dd v
     &\le2\int_0^u(1+2b_0v)^{-2}\dd v
       =\frac{2u}{1+2b_0u}.
 \end{split}
 \label{stat:cost}
\end{equation}
Thus the environmental lower envelope is acquired from the initial
interaction and the real hierarchy. It is not an assumed inequality
throughout the desired time interval.

\begin{corollary}[Three initial gaps]
\label{stat:three-gaps}
Let the target gap at $u=0$ be $g=x_{j+1}-x_j$, and let its two neighboring
gaps be $h_L=x_j-x_{j-1}$ and $h_R=x_{j+2}-x_{j+1}$. Then admissible
initial data for Theorem~\ref{stat:backward-comparison} are
\begin{equation}
 d_0=\frac{g^2}{4},\qquad
 b_0=\frac1{h_L(h_L+g)}+\frac1{h_R(h_R+g)}.
 \label{stat:neighbor-mass}
\end{equation}
Consequently the computable score
\begin{equation}
 R_j(\Delta)=\frac{g_j^2}{4}
               \bigl(1+2b_{0,j}\Delta\bigr)-2\Delta
 \label{stat:score}
\end{equation}
excludes this pair from a first collision through time $\Delta$ whenever
$R_j(\Delta)>0$, subject to the stopped-regime hypotheses.
\end{corollary}

\begin{proof}
For the exterior zero $x_{j-1}$,
$Q=(x_{j-1}-x_j)(x_{j-1}-x_{j+1})=h_L(h_L+g)$.
The right exterior zero contributes $1/[h_R(h_R+g)]$ in the same way.
Every omitted exterior term is positive, so their sum is a lower bound
for $M(0)$. Apply \eqref{stat:strict-margin}.
\end{proof}

One may omit an unavailable neighbor and retain the other positive term.
If only enclosures of the gaps are known, bounds for
\eqref{stat:neighbor-mass} must respect its monotonicities; arbitrary
substitution of enclosure endpoints is not justified. For
$h_L=h_R=g$, we obtain $b_0=1/g^2$, $b_0d_0=1/4$, and
\begin{equation}
 \Delta<\frac{g^2}{6},
 \qquad\text{whereas }b_0=0\text{ gives }\Delta<\frac{g^2}{8}.
 \label{stat:equal-gaps}
\end{equation}
The factor $4/3$ is a local comparison gain. If $g/h_L$ and $g/h_R$
tend to zero, then $b_0d_0\to0$ and the relative gain tends to one.
Thus the equal-gap illustration does not settle the exceptional case of
an isolated narrow gap. Moreover, the smallest initial gap need not
minimize the environment-adjusted score: every eligible pair must be
covered.

Strictness cannot be dropped. For the exact heat polynomial
$H(t,z)=z^2-2(t-t_c)$, one has $D=2(t-t_c)$ and $M=0$.
The comparison is exact and its margin vanishes at
$\Delta=t_*-t_c$, precisely when collision occurs. This finite field
lacks the canonical population but verifies the endpoint interpretation
without numerical evidence.

\subsection{The minimum envelope and its statistical evolution}

Choose $K$ finite buffered blocks. In block $k$, fix a nonempty finite
set $J_k$ of consecutive-pair labels at $u=0$ and follow those labels
throughout the common stopped regime. Boundary-crossing pairs needed
for a covering application must be included; labels are not silently
reselected when zeros move through a window boundary. Set
\begin{equation}
 q_k(u)=\min_{j\in J_k}D_j(u),\qquad
 \mathbb P(X=k)=\frac1K.
 \label{stat:minimum-law}
\end{equation}
Only the block index $X$ is random. The $D_j$ and the canonical divisor
are deterministic.

\begin{proposition}[Weak evolution of the regional minimum]
\label{stat:weak-evolution}
On a compact subinterval of the stopped regime, every $q_k$ is locally
Lipschitz. Choose an active minimizing label $j_k(u)$, with any fixed
rule at ties. For a $C^1$ test function $\varphi$ on the relevant range,
\begin{equation}
 \frac{d}{du}\mathbb E\varphi(q_X(u))
 =\mathbb E\!\left[
    \varphi'(q_X(u))
       \{-2+4q_X(u)M_{j_X(u)}(u)\}\right]
 \quad\text{for almost every }u.
 \label{stat:weak-law}
\end{equation}
Furthermore, throughout that interval,
\begin{equation}
 q_k(u)\ge q_k(0)-2u,
 \qquad
 \mathbb P(q_X(u)\le r)
     \le\mathbb P(q_X(0)\le r+2u),\quad r\ge0.
 \label{stat:tail-transfer}
\end{equation}
\end{proposition}

\begin{proof}
A finite minimum of $C^1$ functions is locally Lipschitz. At a time
where the minimum is differentiable, its derivative equals that of
an active branch; at almost every tie, all tied active derivatives
agree. One can see the latter assertion by applying the elementary
fact that a differentiable function has derivative zero at almost
every point of its zero set to each difference $D_i-D_j$. The finite
union of exceptional null sets remains null. The chain rule and
\eqref{stat:backward-system} therefore give \eqref{stat:weak-law};
the finite average requires no additional interchange assumption.
Positivity gives $D_{j,u}\ge-2$, hence
$D_j(u)\ge D_j(0)-2u$ for every eligible branch. Taking minima proves
the first inequality in \eqref{stat:tail-transfer}. The event
$q_X(u)\le r$ is consequently contained in
$q_X(0)\le r+2u$, which proves the second.
\end{proof}

The minimum may switch branches without invalidating this envelope law.
It does not become one persistent pair solution, and its old pair's
contour, exterior moments and gap anchors do not transfer at a switch.
Equation~\eqref{stat:weak-law} also displays the statistical closure
problem: the minimum marginal depends on the joint law of the minimum
and the complete exterior interaction of its active pair. This is the
relevant BBGKY analogy. No independence or finite marginal closure has
been assumed. Corollary~\ref{stat:three-gaps} identifies a more accessible
initial statistic that bounds a sufficient margin: three consecutive
gaps, with their adjacency conditions retained.

\subsection{Pair counts, exceptional blocks and global coverage}

A one-sided minimum-tail bound does not require the entire minimum law.
At a fixed time let $C_k(r)$ count all unordered distinct-label pairs at
distance at most $r$ in a finite real-zero window. Let $G_k$ be its minimum
gap, with $G_k=+\infty$ when the window has fewer than two occurrences. Then
\begin{equation}
 \{G_k\le r\}=\{C_k(r)\ge1\},\qquad
 \mathbb P(G_X\le r)\le\mathbb E C_X(r).
 \label{stat:pair-majorant}
\end{equation}
The equality follows because a closest pair is consecutive; the inequality
is the pointwise bound $\mathbf1_{\{C_k\ge1\}}\le C_k$.
An exact minimum distribution needs absence-of-points, or void,
information, whereas \eqref{stat:pair-majorant} can be bounded by a
nonnegative two-point majorant. The three-gap score requires more:
four ordered locations must also have no intervening zeros. An
unrestricted four-point marginal cannot silently replace consecutive-gap
data.

The complete marked-pair theory in \cite{BakerOC17,BakerOCReview} gives
time moduli at a fixed positive smoothing scale, retaining multiplicity
and exterior tails. Those results provide neither an initial canonical
profile nor a sharp near-zero gap tail. To use them for
\eqref{stat:pair-majorant}, a smooth nonnegative majorant must control its
resolution, smoothing, boundary and height errors. Their constants grow
as the resolution shrinks, so a sharp-indicator limit is a separate
estimate. Distinct labels at one multiple zero contribute a collision
pair; deleting them as a diagonal would remove the event to be detected.

For the initial score \eqref{stat:score}, define the conservative count
\begin{equation}
 B_\Delta=\#\left\{k:\min_{j\in J_k}R_j(\Delta)\le0\right\}.
 \label{stat:bad-count}
\end{equation}
\begin{proposition}[Finite zero-exception threshold]
\label{stat:integer-exclusion}
If a proved estimate gives
\begin{equation}
 \mathbb P\!\left(\min_{j\in J_X}R_j(\Delta)\le0\right)<\frac1K,
 \label{stat:integer-threshold}
\end{equation}
then every covered pair has a strictly positive score and cannot be a
colliding pair at a first finite collision through backward time
$\Delta$ under the common stopped-regime hypotheses.
\end{proposition}

\begin{proof}
The probability in \eqref{stat:integer-threshold} equals $B_\Delta/K$.
Since $B_\Delta$ is a nonnegative integer, the strict bound implies
$B_\Delta=0$. Corollary~\ref{stat:three-gaps} then applies to every
eligible pair.
\end{proof}

A small probability above this threshold permits exceptions. A continuous
random-center formulation instead needs a proved positive lower measure
of centers that see any bad event, including events near window edges.
Likewise, averaging over observation times can miss an isolated collision;
the event must concern the whole time interval, or carry a proved positive
occupation contribution.

To infer an upper bound $\Lambda\le t_*-\Delta$, these finite comparisons
must be supplemented by a global continuation argument. Its inputs are
a known simple all-real start, coverage of all potentially colliding
pairs including block boundaries, and a uniform high-height exclusion or
complete tail estimate preventing the first failure from escaping to
infinity. An independently proved real/simple high-height statement
uniform over the entire slab could reduce the remaining task to a
buffered compact region. A different cutoff for each fixed time does
not automatically supply this uniform premise. Neither this uniform height
bound nor a canonical estimate strong enough for
\eqref{stat:integer-threshold} is acquired here; density one is not
zero exceptions.

Rodgers and Tao \cite{RodgersTao2020Newman} provide a precedent for using
averaged gap dynamics, but with a different conclusion. Under the
contradictory assumption $\Lambda<0$, their nearly equal spacing
consequence conflicts with a positive proportion of small gaps, proving
$\Lambda\ge0$. Their negative-time motion theorem is not used to justify
our positive-time identity; Lemma~\ref{stat:strict-isolation} supplies
that passage on \eqref{stat:slab}. Here the proposed application is a
lower \emph{upper bound}, obtained by uniform backward preservation of
reality. The established gain is the explicit local comparison and its
statistical consumer, not a new numerical bound on $\Lambda$.
